\documentclass[11pt]{article}
\usepackage{booktabs}

\begin{document}

\section{Surprisal at the disambiguating region}

We score every item with both models and take the mean log-probability at the
disambiguating word. Table~\ref{tab:surprisal} reports the result: both models
assign higher surprisal to the ambiguous condition, and GPT-2-large shows the
larger effect.

\begin{table}[t]
\centering
\begin{tabular}{lcc}
\toprule
Model & Ambiguous & Unambiguous \\
\midrule
GPT-2-large & 8.42 & 6.10 \\
Pythia-1.4B & 7.95 & 6.33 \\
\bottomrule
\end{tabular}
\caption{Mean surprisal in bits at the disambiguating region, NP/Z items.}
\label{tab:surprisal}
\end{table}

The gap between conditions is what we call the garden-path effect throughout.

\end{document}
